\par\medskip
\noindent
\underline{Teorema} 4.3. \underline{Sea $S$ un rasgo henseliano de
cuerpo residual $k$ y cuerpo de funciones $K$, y sea $A_K$ un esquema
abeliano sobre $K$. Considérese la representación de}
$\pi=\operatorname{Gal}(\overline K/K)$ \underline{sobre el módulo de
Tate} $U_\ell=T_\ell(A(\overline K))$, \underline{donde $\ell$ es un
primo distinto de la característica residual.}

\smallskip
\noindent
a) \underline{Existe un subgrupo abierto $I'$ del grupo de inercia que
actúa unipotentemente sobre $U_\ell$, de modo que el carácter de la
representación de $\pi$ sobre $U_\ell$ es, de hecho, una función sobre
$\pi/I'$.}
\underline{Además, sobre $\Gamma=I/I'$ esta función toma valores enteros
y es independiente de $\ell$. Más generalmente, para todo $g\in I$,
los coeficientes del polinomio característico de $g_{U_\ell}$ son
enteros racionales independientes de $\ell$.}

\smallskip
\noindent
b) \underline{Supongamos que $k$ es un cuerpo finito de $q$ elementos,
de modo que}
\[
 \pi_o=\pi/I\simeq\operatorname{Gal}(\overline k/k)
\]
\underline{se identifica con el completado profinito
$\widehat{\mathbb Z}$ de $\mathbb Z$, correspondiendo al generador $1$
de $\mathbb Z$ el automorfismo de Frobenius}
\[
 f_q:x\longmapsto x^q
\]
\underline{de $\overline k/k$. Entonces, para todo $g\in\pi$ cuya imagen
en $\pi_o$ pertenece a $\mathbb Z$, la traza} (\underline{y, más generalmente, todo
coeficiente del polinomio característico}) \underline{de $g_{U_\ell}$ es un número
racional independiente de $\ell$. Es incluso un entero si la imagen de
$g$ en $\pi_o$ es de la forma $f_q^{\,n}$, con $n\geq0$.}

Demostremos, por ejemplo, las afirmaciones relativas a las trazas. Las
afirmaciones relativas a los coeficientes del polinomio característico
podrían demostrarse del mismo modo. Puede observarse también que se
deducen formalmente de las afirmaciones sobre las trazas, expresando
estos coeficientes en términos de las trazas de los iterados (aunque
estas fórmulas introducen denominadores y, por tanto, no dan la
integralidad afirmada), junto con un lema de Dwork--Fatou, del que
Kleiman da una elegante demostración en [14, 2.8].

\smallskip
\noindent
a) Utilizamos la fórmula (4.2.10), que aquí se convierte en
\begin{equation}\tag{4.3.0}
 \operatorname{Tr}g_{U_\ell}
 =\operatorname{Tr}g_{B^*_o,\ell}
 +\operatorname{Tr}g_{T^*_o,\ell}
 +\operatorname{Tr}g^{-1}_{T^*_o,\ell}.
\end{equation}
Aquí $g_{B^*_o}$ y $g_{T^*_o}$ son ahora $k'$-automorfismos de $B^*_o$
y $T^*_o$, respectivamente. Basta, por tanto, demostrar que las dos
trazas correspondientes son enteros racionales independientes de
$\ell$, lo cual es una afirmación «geométrica» general sobre grupos
algebraicos. Para $g_{B^*_o}$ es un resultado bien conocido de Weil en
la teoría de las variedades abelianas [5] [15]. Para $g_{T^*_o}$ es
prácticamente inmediato: se describe la representación de $\Gamma$
sobre $T=T^*_o$ mediante la representación correspondiente sobre el
grupo dual $M'\simeq\mathbb Z^\mu$. La traza buscada es la de la
contragrediente de esta última representación, por el isomorfismo
canónico
\[
 T_\ell(T)\simeq\check M_\ell(1)
 \mathrel{\mathop{=}^{\mathrm{def}}}
 \check M'\otimes_{\mathbb Z}\mathbb Z_\ell(1).
\]

\smallskip
\noindent
b) Si la imagen de $g$ en $\pi_o$ es $f_q^{\,n}$, la fórmula (4.2.10)
se convierte ahora en
\begin{equation}\tag{4.3.1}
 \operatorname{Tr}g_{U,\ell}
 =\operatorname{Tr}g_{B^*_o,\ell}
 +\operatorname{Tr}g_{T^*_o,\ell}
 +q^n\operatorname{Tr}g^{-1}_{T^*_o,\ell}.
\end{equation}
Quedamos reducidos a demostrar el lema siguiente, en el que $k$ debe
sustituirse por una clausura algebraica de $k'$, y $A$ por
$A^{*o}_o\otimes_{k'}\overline k$.

\medskip
\noindent
\underline{Lema} 4.3.2. \underline{Sea $k$ la clausura algebraica de
un cuerpo finito, sea $A$ un esquema en grupos conmutativo de tipo
finito sobre $k$, y sea $g$ un automorfismo de este esquema en grupos
compatible con el automorfismo $f$ de $\operatorname{Spec}(k)$ obtenido
por transporte de estructura a partir de un automorfismo de Frobenius
iterado}
\[
 x\longmapsto x^q,
 \qquad q=p^r,\quad r\geq0,
\]
\underline{donde $q$ es una potencia de la característica. Así, la
acción de $g$ sobre las funciones $\varphi$ satisface}
\[
 g^*(\lambda\varphi)=f^*(\lambda)g^*(\varphi)
 =\lambda^{q^{-1}}g^*(\varphi)
 =\lambda^{p^{-r}}g^*(\varphi)
 \qquad(\lambda\in k).
\]
\underline{Sea $g_\ell$ el automorfismo inducido por $g$ sobre
$T_\ell(A(k))$ por «transporte de estructura», donde $\ell$ es un primo
distinto de $p$. Entonces la traza de $g_\ell$ (y, más generalmente,
todo coeficiente de su polinomio característico) es un entero racional
independiente de $\ell$, y sus valores propios son enteros algebraicos
de valor absoluto $q^{1/2}$ o $q$. Si $A$ es un esquema abeliano, sus
valores absolutos son $q^{1/2}$; si $A$ es afín, los valores propios son
de la forma $q\varepsilon_i$, donde los $\varepsilon_i$ son raíces de
la unidad.}

En efecto, considérese el endomorfismo «potencia $q$» $f_q^A$ del
esquema absoluto subyacente de $A$, definido como la identidad sobre
los puntos y como la aplicación potencia $q$ sobre el haz estructural.
Es un endomorfismo del esquema en grupos $A$ compatible con el
endomorfismo $f_q^{\operatorname{Spec}(k)}=f^{-1}$ de la base
$\operatorname{Spec}(k)$. En consecuencia, el compuesto
\begin{equation}\tag{4.3.3}
 F=f_q^A g
\end{equation}
es un endomorfismo del esquema en grupos $A$ compatible con la identidad
sobre $\operatorname{Spec}(k)$; en otras palabras, es un endomorfismo
sobre $k$ del esquema en grupos $A$. Afirmo que el automorfismo
$\Gamma_g$ de $A(k)$ definido por $g$ mediante transporte de estructura,
a saber,
\begin{equation}\tag{4.3.4}
 \Gamma_g(u)=g\circ u\circ f^{-1}
 \qquad\text{para }u\in A(k),
\end{equation}
puede describirse también por
\begin{equation}\tag{4.3.5}
 \Gamma_g(u)=F(u)
 \qquad\text{para }u\in A(k).
\end{equation}
Para comprobarlo, basta observar que las aplicaciones subyacentes de
conjuntos de los morfismos $\operatorname{Spec}(k)\longrightarrow A$
representados por ambos miembros son iguales. Esto nos reduce a la
igualdad obtenida de (4.3.5) componiendo su miembro izquierdo por la
izquierda con $f_q^A$ y su miembro derecho por la derecha con $f^{-1}$,
puesto que $f_q^A$ y $f^{-1}$ inducen la identidad sobre los conjuntos
subyacentes de los esquemas fuente. La fórmula resultante se deduce
inmediatamente de (4.3.3) y (4.3.4).

Aplicando (4.3.5) a $u\in{}_{\ell^\nu}A(k)$ y pasando al límite inverso
en $\nu$, se obtiene que $g_\ell$ puede interpretarse también como el
endomorfismo $T_\ell(F)$ de $T_\ell(A(k))$ asociado al endomorfismo $F$
de $A$. Por otra parte, se sabe que $T_\ell(F)$ posee un polinomio
característico con coeficientes enteros independientes de $\ell$. Por
dévissage, esto se reduce al caso en que $A$ es un esquema abeliano o
un toro, donde el resultado requerido ya fue utilizado en la parte a)
de la demostración. El mismo argumento demuestra también las
afirmaciones sobre los valores propios, en vista, en el caso de los
esquemas abelianos, de los resultados clásicos de Weil [15].

Esto demuestra 4.3.2 y concluye la demostración de 4.3. Al mismo tiempo,
demuestra el resultado siguiente.

\medskip
\noindent
\underline{Corolario} 4.4. \underline{Supongamos que se cumplen las
hipótesis de 4.3 b). Entonces, para todo $g\in\pi'$ cuya imagen en
$\pi_o$ es de la forma $f_q^n$ con $n\geq0$, los valores propios de los
automorfismos inducidos de}
\[
 U^{\mathrm{et}},\qquad
 U^{\mathrm{ef}}/U^{\mathrm{et}},\qquad
 U/U^{\mathrm{ef}}
\]
\underline{son enteros algebraicos cuyos valores absolutos son,
respectivamente, $q^n$, $q^{n/2}$ y} $1$. \underline{Para los dos trozos extremos
puede decirse más precisamente que los valores propios correspondientes
son, respectivamente,}
\[
 (q^n\varepsilon_i)_{1\leq i\leq\mu}
 \qquad\text{y}\qquad
 (\varepsilon_i)_{1\leq i\leq\mu},
\]
\underline{donde los $\varepsilon_i$ $(1\leq i\leq\mu)$ son raíces de
la unidad.}

En efecto, la fórmula (4.1.2) da para el último trozo los valores propios
\[
 \frac{q^n}{q^n\varepsilon_i}=\varepsilon_i^{-1}.
\]
Pero, puesto que $\varepsilon_i$ es una raíz de la unidad, se tiene
$\varepsilon_i^{-1}=\overline{\varepsilon_i}$ (conjugado complejo), y
como la familia de los $\varepsilon_i$ es estable bajo conjugación
compleja (por ser la familia de raíces de una ecuación con coeficientes
enteros), la familia de los $\varepsilon_i^{-1}$ es también la familia
de los $\varepsilon_i$. Esto demuestra la última afirmación de 4.4.
